\section{Macroscopic crossings}\label{sec:crossings}

This section proves the dichotomy for arbitrary linear-time traces (Theorem~D) and its consequence for oblivious transposers
(Theorem~X). Both rest on one implication, proved for every trace without any hypothesis on what the machine computes: if a
small set of steps cuts a large set of outputs off from the crossing steps $\CL_H$, then the remaining cone relation contains a
linear-size acyclic set. \Cref{sec:certificates} has shown that such an acyclic set is incompatible with transposition.

\subsection{Constants}\label{sec:constants}

Fix $k\ge1$, reals $C\ge1$ and $\alpha>0$, and a depth $D\ge0$. Put
\[
  m_0(k)=(2k+1)(k+1)^{2k},\qquad K_{\mathrm{row}}=\lceil24C/\alpha\rceil,\qquad B_{\mathrm{time}}=\lceil8m_0C/\alpha\rceil,
\]
\[
  \rho_0(k,C,\alpha)=\frac{1}{3\binom{m_0+2K_{\mathrm{row}}}{m_0}},\qquad
  N_{\mathrm{thr}}=(2k+1)+m_0(2k+1)(B_{\mathrm{time}}+1),\qquad
  E_0(k,D,C,\alpha)=(k+2)^D N_{\mathrm{thr}}(5k+1),
\]
\[
  \theta_{\mathrm{E}}(k,D,C,\alpha)=\frac{\alpha}{12\,E_0(k,D,C,\alpha)},\qquad
  d_{\mathrm{pb}}(k,C,\alpha)=\max\Bigl(K_{\mathrm{row}}B_{\mathrm{time}},\Bigl\lceil\frac{12m_0(2B_{\mathrm{time}}+3)}{\alpha}\Bigr\rceil\Bigr).
\]
Here $m_0(k)$ counts the \emph{phase-exit chains} of the proof below; $m_0=12$, $405$ and $28672$ for $k=1,2,3$. The constant
$\rho_0$ does not depend on $D$ or on the alphabet. For a machine $M$ with $k$ work tapes and value width $\wval$, and $\alpha>0$,
put
\[
  \cE(M,C,\alpha)=\frac{\rho_0(k,C,\alpha)}{2\wval},\qquad
  D_{\mathrm{E}}(M,C,\alpha)=\mathrm{genieDepth}\bigl(k,\wval,C,\alpha\rho_0(k,C,\alpha)/4\bigr),
\]
\[
  \thetaM(M,C,\alpha)=\theta_{\mathrm{E}}\bigl(k,D_{\mathrm{E}}(M,C,\alpha/2),C,\alpha/2\bigr)
  =\frac{\alpha}{24\,E_0\bigl(k,D_{\mathrm{E}}(M,C,\alpha/2),C,\alpha/2\bigr)},
\]
where $\mathrm{genieDepth}$ is the depth bound of the fourteenth paper's genie (\cref{lem:genie}). In Lean these are
\lean{conjectureE_m0}, \lean{conjectureE_bridge_rho0} (through \lean{timePullbackRhoMin}), \lean{conjectureE_Nthresh},
\lean{conjectureE_E0}, \lean{conjectureE_theta}, \lean{timePullback_gen_d0}, \lean{conjectureE_cE}, \lean{conjectureE_DE} and
\lean{conjectureE_thetaMacro}.

These constants are positive and astronomically small. For $k=2$, $C=1$: $1/\rho_0(2,1,1)\approx10^{65.8}$ and
$1/\rho_0(2,1,\tfrac12)\approx10^{105.4}$; for $k=1$ and $k=3$ at $C=\alpha=1$ the values are $10^{12.6}$ and $10^{153.4}$. The
threshold $\thetaM$ is far smaller, because $E_0$ grows like $(k+2)^{D}$ and the genie depth $D_{\mathrm{E}}$ is itself at least
$32\wval C/(\alpha\rho_0)$: for $k=2$ and $C=\alpha=1$, $\thetaM\le4^{-10^{107}}$. Nothing in the problem suggests that the true
values are this small (\cref{rem:numerics-crossings}).

\subsection{The time-pullback bound}

The first ingredient turns a structure in time into an acyclic set. Split $\Range T$ into $B_{\mathrm{time}}$ consecutive
\emph{time blocks} $T_0,\dots,T_{B_{\mathrm{time}}-1}$ of length $\lfloor T/B_{\mathrm{time}}\rfloor$, the last one extended to end
at~$T$.

\begin{theorem}[Time pullback]\label{thm:pullback}
\formal{MacroscopicThreshold.lean}{theorem_51_2_general_zbad_mais}
Let $\tau$ be a trace with $k$ work tapes that appends exactly $L=d^2b$ bits before step $T\le C(L+1)$, with $d,b,m_0\ge1$, $C\ge1$ and
$\alpha>0$, and let $d\ge d_{\mathrm{pb}}$ computed with this $m_0$. Let $V_0\subseteq\Range L$ with $\abs{V_0}\ge\alpha L$, let
$Z\subseteq\Range L$ with $\abs Z\le\abs{V_0}/12$, let $B$ be a set of steps and $D\ge0$. Suppose there are at most $m_0$ indices
$\nu$, each with a nondecreasing sequence $\xi_\nu(0)\le\dots\le\xi_\nu(B_{\mathrm{time}})$ and a permutation $\pi_\nu$ of the
blocks, such that the following holds for every $q\in V_0\setminus Z$: whenever a dependency path from $\Out(q)$ with at most $D$
edges, whose vertices after the root avoid $B$, ends at a step of block $T_\beta$ at which the input head reads a position of $V_0$,
the root $\Out(q)$ lies in $[\xi_\nu(\pi_\nu(\beta)),\xi_\nu(\pi_\nu(\beta)+1))$ for some~$\nu$. Then the ball relation of radius
$D$ avoiding $B$ contains an acyclic set $U\subseteq V_0$ with
\[
  \abs U\ \ge\ \frac{\abs{V_0}}{3\binom{m_0+2K_{\mathrm{row}}}{m_0}} .
\]
\end{theorem}

\begin{outline}
For each $q$ let $A(q)$ be the set of blocks $\beta$ such that $\Out(q)$ lies in an interval assigned to $\beta$; since the
intervals of one $\nu$ are disjoint, $\abs{A(q)}\le m_0$. For an input position $p$ let $E(p)$ be the set of blocks in which $p$
is read. The hypothesis says that if $p$ is in the ball of $q$ then $A(q)$ and $E(p)$ meet. A kinematic count
(\lean{kinematic_goodInputRows_card_ge}) shows that all but at most $3T/(K_{\mathrm{row}}W)$ input rows are touched by the input head
in at most $K_{\mathrm{row}}$ runs, each shorter than a block, so $\abs{E(p)}\le2K_{\mathrm{row}}$ off a quarter of $V_0$. Pairs
with $A(q)\cap E(q)\ne\emptyset$, self-collisions, are bounded without any anti-concentration assumption, using the row and column
structure of $\tpos$ and the monotonicity of append steps: for a fixed $\nu$ each threshold splits at most one output row
(\lean{card_selfCollisions_singleFamily_le}). Off the bad sets the relation is contained in the \emph{block relation}
$A(q)\cap E(p)\ne\emptyset$ with $\abs A\le m_0$, $\abs E\le2K_{\mathrm{row}}$ and no self-loops, and a uniform block cut
(\lean{exists_isAcyclicIn_card_ge_of_uniformBlockCut}) extracts an acyclic set of the stated size.
\end{outline}

\Cref{thm:pullback} is purely about the trace. The constant $1/\binom{m_0+2K_{\mathrm{row}}}{m_0}$ is the price of the uniform
block cut over all possible block sets, and it is the source of the size of $\rho_0$.

\subsection{The bridge}

The second ingredient finds the time-pullback structure whenever $\CL_H$ can be cut off. We need a separation notion that
ignores the roots themselves.

\begin{definition}[Separation within depth $D$]\label{def:ballsep}
\formal{EnvelopeCrux.lean}{IsBallPathSeparator}
Let $V_0$ be a set of input positions and $S$, $B$ sets of steps. $B$ \emph{separates $V_0$ from $S$ within depth $D$} if every
dependency path from a root $\Out(q)$, $q\in V_0$, with at most $D$ edges and with all its vertices after the root outside $B$,
also has all its vertices after the root outside $S$ (\lean{IsBallPathSeparator}).
\end{definition}

If $B$ separates $\Out(V_0)$ from $S$ in the sense of \cref{def:cut}, then $B$ separates $V_0'=\{q\in V_0:\Out(q)\notin B\}$
from $S$ within every depth (\lean{isBallPathSeparator_of_isMengerVertexCut}).

\begin{theorem}[Bridge]\label{thm:bridge}
\formal{MacroscopicThreshold.lean}{conjectureE_bridge_exists_isAcyclicIn_crossingLowerGT2H, conjectureE_bridge_exists_isAcyclicIn_crossingLowerH}
Let $\tau$ be a trace with $k$ work tapes that appends exactly $L=d^2b$ bits before step $T\le C(L+1)$, with $C\ge1$ and
$\alpha>0$. Let $D\ge0$, let $1\le H\le\theta_{\mathrm{E}}(k,D,C,\alpha)\,L$, and let $d\ge d_{\mathrm{pb}}(k,C,\alpha)$. If $V_0$ has
at least $\alpha L$ elements and $B$ separates $V_0$ from $\CL_H$ within depth $D$, then the ball relation of radius $D$ avoiding
$B$ contains an acyclic set $U\subseteq V_0$ with $\abs U\ge\rho_0(k,C,\alpha)\abs{V_0}$.
\end{theorem}

\begin{outline}
\emph{Two cuts per tape.} On every unit-speed walk there are two steps $\vartheta_+,\vartheta_-$ such that every revisit of
excursion greater than $H$ that crosses no other such revisit straddles one of them: the upper ones straddle $\vartheta_+$, the
lower ones $\vartheta_-$ (\lean{exists_two_cuts_covering_nonCrossingLargeHops}). \emph{Dichotomy along paths.} Along a dependency
path from $V_0$ that avoids $B$ after the root, every edge is a state edge, a revisit edge of excursion at most $2H$, or a revisit
edge whose revisit crosses nothing large and therefore straddles one of the $2k$ cut steps; the separation hypothesis excludes
large crossing revisits (\lean{depPath_nonCrossing_dichotomy_of_isBallPathSeparator_crossingLowerH}). \emph{Phases and chains.}
The cut steps split time into at most $2k+1$ work phases; a path visits the phases along one of at most $m_0(k)$ phase-exit chains
(\lean{card_chainIdx_le}), and along a chain each passage between phases is a strictly antitone partial bijection between times
(\lean{crossingCut_nonCrossing_strictAnti}). Pulling the block boundaries back along each chain gives, for every chain, monotone
root intervals of both parities, which is the structure required by \cref{thm:pullback} with $m_0=m_0(k)$. \emph{Guards.} Small
revisits and the state edge are handled by a guard set of at most $(5k+1)H$ steps per threshold
(\lean{card_microHopStepGuard_le}); the roots that reach a guard within $D$ edges form the exceptional set $Z$, of size at most
$(k+2)^D N_{\mathrm{thr}}(5k+1)H=E_0H\le\alpha L/12$ by the choice of~$H$ (\lean{card_ancestors_le}). \Cref{thm:pullback} applies.
\end{outline}

The second Lean declaration is the same statement for banded crossing sets with an explicit clause on overshooting revisits;
for the unbanded set $\CL_H$ the clause holds for every trace (\lean{clauseE1H_uncapped}). The form with $H=W$ is
\lean{conjectureE_bridge_exists_isAcyclicIn_crossingLower}.

\subsection{Theorem D: a dichotomy for every linear-time trace}

\begin{theorem}[Theorem~D]\label{thm:D}
\formal{MacroscopicThreshold.lean}{linearTimeTrace_crossingLowerGT2H_menger_dichotomy, linearTimeTrace_crossingLowerGT2H_ballSep_dichotomy}
Let $\tau$ be any trace with $k$ work tapes that appends exactly $L=d^2b$ bits before step $T\le C(L+1)$, with $C\ge1$ and
$\alpha>0$. Let $D\ge0$, $1\le H\le\theta_{\mathrm{E}}(k,D,C,\alpha/2)\,L$, $d\ge d_{\mathrm{pb}}(k,C,\alpha/2)$, and $c\le\frac12$.
For every set $V_0$ of at least $\alpha L$ input positions, exactly one of the following holds:
\begin{enumerate}[label=(\roman*)]
\item some set $B$ of fewer than $c\abs{V_0}$ steps separates $\Out(V_0)$ from $\CL_H$, and the ball relation of radius $D$ avoiding
$B$ contains an acyclic set $U\subseteq V_0$ with $\abs U\ge\frac12\rho_0(k,C,\alpha/2)\abs{V_0}$;
\item every set separating $\Out(V_0)$ from $\CL_H$ has at least $c\abs{V_0}$ elements.
\end{enumerate}
\end{theorem}

\begin{outline}
If (ii) fails, take a separating set $B$ with $\abs B<c\abs{V_0}\le\abs{V_0}/2$. Since $\Out$ is injective, at least half of the
roots of $V_0$ lie outside $B$; the set $V_0'$ of these has at least $(\alpha/2)L$ elements, and $B$ separates it from $\CL_H$
within depth~$D$. \Cref{thm:bridge} at density $\alpha/2$ gives $U$.
\end{outline}

Theorem~D mentions neither a machine nor transposition. Its content is the implication from (i)'s first half to its second: a
cheap cut from the crossing steps leaves a large acyclic set. Read the other way, it says that a trace whose cone relations have
no large acyclic set after cheap cuts must reach $\CL_H$ with linear expansion from every large output set. The transposition
hypothesis enters only through the certificate method.

\subsection{Theorem X}

\begin{theorem}[Theorem~X, macroscopic crossings]\label{thm:X}
\formal{MacroscopicThreshold.lean}{conjectureE_obliviousLinTime_macroscopicCrossingLowers_menger_lb, conjectureE_obliviousLinTime_crossingLowerGT2H_menger_lb}
Let $M$ be a machine with $k$ work tapes and value width $\wval$, and let $C\ge1$ and $\alpha>0$. Put $\theta=\thetaM(M,C,\alpha)$ and
\[
  d_0=\max\Bigl(1,\ d_{\mathrm{pb}}(k,C,\alpha/2),\ \Bigl\lceil\frac{24E_0(k,D_{\mathrm{E}}(M,C,\alpha/2),C,\alpha/2)}{\alpha}\Bigr\rceil,\
  \Bigl\lceil\frac2\theta\Bigr\rceil\Bigr).
\]
Let $d\ge d_0$, suppose $\mathrm{ObliviousLinTime}(M,d,C,T)$, and let $H^*=\lfloor\theta L\rfloor$. Then $W\le H^*$ and
$\theta L\le2H^*$, and for every set $V_0$ of at least $\alpha L$ input positions, every set of steps separating $\Out(V_0)$ from
$\CL_{H^*}$ in the dependency DAG of the common trace has at least
\[
  \tfrac12\,\cE(M,C,\alpha/2)\,\abs{V_0}\ =\ \frac{\rho_0(k,C,\alpha/2)}{4\wval}\,\abs{V_0}
\]
elements. The same bound holds for $\CL_H$ at every $1\le H\le\theta L$ and every $d\ge\max(1,d_{\mathrm{pb}}(k,C,\alpha/2))$.
\end{theorem}

Every revisit counted in $\CL_{H^*}$ has excursion greater than $2H^*\ge\theta L$: these are crossings of macroscopic revisits.
When $\alpha\le1$, taking $V_0$ to be a set of $\lceil\alpha L\rceil$ input positions and the separating set to be $\CL_{H^*}$ itself gives the
following.

\begin{corollary}[Linearly many macroscopic crossings]\label{cor:many}
\formal{MacroscopicThreshold.lean}{conjectureE_obliviousLinTime_macroscopicCrossingLowers_menger_lb}
Under the hypotheses of \cref{thm:X}, if moreover $\alpha\le1$, the common trace has at least $\frac{\alpha\rho_0(k,C,\alpha/2)}{4\wval}L$ steps that are
earlier ends of crossing revisits of excursion greater than $\theta L$.
\end{corollary}

\begin{outline}[Outline of the formal proof of \cref{thm:X,cor:many}]
\emph{Thresholds.} From $d\ge d_0$ one gets $W\le\theta L$, so $1\le W\le H^*$, and $\theta L\ge2$, so $\theta L<H^*+1\le2H^*$.
\emph{Reduction.} Let $B$ separate $\Out(V_0)$ from $\CL_{H^*}$. If $\abs B\ge\abs{V_0}/2$ we are done, since
$\frac12\cE\le\frac1{24}$. Otherwise the set $V_0'$ of $q\in V_0$ with $\Out(q)\notin B$ has at least $\abs{V_0}/2\ge(\alpha/2)L$
elements, and $B$ separates it from $\CL_{H^*}$ within the depth $D_{\mathrm{E}}(M,C,\alpha/2)$. \emph{Bridge.} By the definition of
$\thetaM$, $H^*\le\theta_{\mathrm{E}}(k,D_{\mathrm{E}},C,\alpha/2)L$, and \cref{thm:bridge} at density $\alpha/2$ gives an acyclic
$U\subseteq V_0'$ in the ball relation with $\abs U\ge\rho_0\abs{V_0'}$. \emph{Coding.} \Cref{lem:acyclic-cost} with the genie of
cost $(\alpha\rho_0/8)L$ and depth $D_{\mathrm{E}}$ gives $\wval\abs B\ge\abs U-(\alpha\rho_0/8)L\ge(\rho_0/2)\abs{V_0'}$, hence
$\abs B\ge\cE\abs{V_0'}\ge\frac12\cE\abs{V_0}$.
\end{outline}

\begin{remark}[The earlier, row-scale form]\label{rem:2W}
With $H=W$ the threshold condition is automatic for large $d$, and the conclusion holds with $\CL_W$, the crossing steps of
revisits longer than two rows (\lean{conjectureE_obliviousLinTime_crossingLowerGT2W_menger_lb}); the same bound holds with the
larger set of both ends of revisits that cross or share a cell and a turning point with another revisit
(\lean{conjectureE_obliviousLinTime_mixGT2W_menger_lb}). Theorem~X is stronger: raising the threshold from $2W$ to $\theta L$
shrinks the target set, so rearrangements below a fixed fraction of $L$ can never substitute for macroscopic crossings.
\end{remark}

\begin{remark}[Numerical]\label{rem:numerics-crossings}
We computed exact minimum separating sets, by maximum flow, from large output sets to the earlier ends of crossing revisits of
excursion greater than $E$, for $E\in\{2W,L/64,L/16,L/4,L/2,L\}$, on several explicit oblivious schedules at
$d\in\{32,64,128\}$ and $b=1$. On the five linear-time cascade schedules of \cref{sec:cascades}, and on two superlinear transposers,
the ratio of the minimum separating set to $\abs{V_0}$ was exactly $1$ for every $E\le L/2$. At $E=L$ three of the cascades make no
such crossings at all, so the empty set separates, while heuristic bounds on their largest acyclic sets fell towards~$0$ as the
radius grew. So, numerically, the threshold in Theorem~D appears not to be a universal constant such as $\frac12$. These measurements are far below the
sizes at which the theorems apply and are no evidence for them; they suggest only that the true expansion constant may be of
order~$1$ at thresholds of order $L/C$.
\end{remark}
